\subsection{Fenwick Tree con XOR + Sprague-Grundy (Nim por rangos)}

\graybold{Preguntas guía:}

\begin{itemize}
\item ¿El juego es una suma de sub-juegos independientes (aquí: columnas de fragmentos) donde en cada turno se modifica UNA sola pila?
\item ¿Necesito responder "quién gana" para muchos rangos distintos, con actualizaciones puntuales intercaladas?
\item ¿La operación de combinación relevante es XOR en vez de suma?
\end{itemize}

\graybold{Utilidad:} Generalizar el Fenwick Tree (visto para sumas) a CUALQUIER operación invertible para mantener el valor de Grundy combinado de un rango de pilas de Nim, con actualizaciones puntuales en tiempo logarítmico. Combinado con el teorema de Sprague-Grundy/Bouton, responde instantáneamente quién gana un juego de Nim sobre cualquier subrango.

\graybold{Complejidad:} Construcción \bigO{n}. Update puntual \bigO{\log n}. Query de rango \bigO{\log n}.

\graybold{Fundamento teórico:} En Nim con pilas independientes donde un jugador retira cualquier cantidad positiva de UNA pila, el valor de Grundy de una pila de tamaño $v$ es exactamente $v$. Por el teorema de Bouton, el jugador que mueve PRIMERO gana si y solo si el XOR de todas las pilas en juego es distinto de 0. Como XOR es su propio inverso ($v \oplus v = 0$), el mismo esquema del Fenwick Tree que acumula sumas de prefijo funciona igual acumulando XOR de prefijo: para modificar el valor de una posición de $v_{viejo}$ a $v_{nuevo}$, basta aplicar un update con delta $= v_{viejo} \oplus v_{nuevo}$ (equivalente a "quitar" el valor viejo y "poner" el nuevo en una sola operación).

\graybold{Ejercicio aplicado}

Juan y Frank juegan Nim sobre columnas de fragmentos: en cada turno, un jugador retira cualquier cantidad positiva de fragmentos de UNA columna; pierde quien no puede mover. Dadas Q consultas, cada una una "profecía" (determinar el ganador en un rango $[l,r]$, sabiendo que Frank siempre empieza) o un "ritual" (sumar $x$ fragmentos a una columna), resolver cada profecía.

\graybold{I/O:} N,Q ≤ \ten{6} → lectura total con tokenización (patrón 2), y las funciones de update/query del Fenwick se escriben SIN llamadas a función (inline) para reducir el overhead de Python en este volumen de operaciones. Aun así, en el peor caso (N=Q=\ten{6}) la solución pura en Python ronda \aprox{}7s si el juez es muy estricto con el límite de tiempo, este es el techo práctico de Python puro para este volumen sin cambiar de lenguaje.

\graybold{Código solución}

\begin{lstlisting}[language=Python]
import sys

def main():
    data = sys.stdin.buffer.read().split()
    idx = 0
    n = int(data[idx]); q = int(data[idx + 1]); idx += 2
    arr = [0] * (n + 1)
    tree = [0] * (n + 1)

    # construccion en O(n): cada nodo se combina UNA vez con su padre directo,
    # en vez de hacer n updates completos de O(log n) cada uno
    for i in range(1, n + 1):
        v = int(data[idx]); idx += 1
        arr[i] = v
        tree[i] ^= v
        j = i + (i & (-i))
        if j <= n:
            tree[j] ^= tree[i]

    out = []
    out_append = out.append
    ANS = ("JUAN", "FRANK")
    P_BYTE = 80  # ord('P'): comparar el primer byte evita decodificar el token

    i2 = idx
    for _ in range(q):
        tok = data[i2]
        if tok[0] == P_BYTE:
            l = int(data[i2 + 1]); r = int(data[i2 + 2]); i2 += 3
            # xor(prefijo r) ^ xor(prefijo l-1) = xor del rango [l, r]
            res = 0
            i = r
            while i:
                res ^= tree[i]
                i &= i - 1          # equivalente a i -= i & (-i), un poco mas rapido
            i = l - 1
            while i:
                res ^= tree[i]
                i &= i - 1
            out_append(ANS[res != 0])   # XOR != 0 -> gana quien mueve primero (Frank)
        else:
            k = int(data[i2 + 1]); x = int(data[i2 + 2]); i2 += 3
            old = arr[k]
            new = old + x
            delta = old ^ new           # "sacar" old y "poner" new en un solo update
            arr[k] = new
            i = k
            while i <= n:
                tree[i] ^= delta
                i += i & (-i)

    sys.stdout.write("\n".join(out) + "\n")

main()
\end{lstlisting}
